% =====================================================================
%  Widma EDS próbek o różnym składzie (pomiary 13.11.2026) — dalsza część
%  sekcji \ref{sec:widma}. Wstawiany we wspolne/rozdzialy/05_badania.tex
%  po % =====================================================================
%  Widma EDS — fragment rozdziału 5 (sekcja \ref{sec:widma}).
%  Wstawiany we wspolne/rozdzialy/05_badania.tex po \section{Widma EDS próbek}.
%  Dane: widma-eds/dane_mca/2026-06-12_Al_4kV/4SERIA.mca (wykres: skrypty_mca/widma_mca.py).
%  Rysunki: widma-eds/widmo_eds_4kV.pdf, widma-eds/stabilnosc_zliczen.pdf
% =====================================================================

\subsection{Widmo próbki aluminiowej przy napięciu 4~kV}
\label{sec:widma-warunki}

Widmo zarejestrowano 12.06.2026 przy napięciu przyspieszającym
$\UT=\SI{4000}{\volt}$ i~napięciu ekstrakcyjnym $\UB=\SI{1000}{\volt}$
spektrometrem opisanym w~\pref{sec:rejestracja-widm}. Ustawienia toru
akwizycji zestawiono w~\tabref{tab:widma-parametry}.

\begin{table}[htbp]\centering
  \caption{Parametry rejestracji widma EDS}\label{tab:widma-parametry}
  \begin{tabular}{@{}ll@{}}\toprule
    Parametr & Wartość \\\midrule
    Napięcie przyspieszające $\UT$ & \SI{4000}{\volt} \\
    Napięcie ekstrakcyjne $\UB$ & \SI{1000}{\volt} \\
    Czas kształtowania (\angielski{peaking time}) & \SI{0.5}{\micro\second} \\
    Liczba kanałów, szerokość kanału & 8192, \SI{10}{\electronvolt} \\
    Czas rzeczywisty / czas życia & \SI{210.7}{\second} / \SI{100.1}{\second} \\
    Częstość zliczeń wejściowa / wyjściowa & ok.~475~kcps / ok.~225~kcps \\
    Czas martwy & \SIrange{52}{54}{\percent} \\\bottomrule
  \end{tabular}
\end{table}

Widmo zbierano w~trybie ciągłym z~kumulacją, a~stan pomiaru zapisano
w~ośmiu chwilach pośrednich. Badaną próbką była płytka aluminiowa, a~pomiar
wykonano z~magnesem ogniskującym umieszczonym na emiterze. Ciśnienie
w~komorze nie przekraczało \SI{1e-5}{\milli\bar}.
\TODO{prąd emisji $\Ie$ podczas pomiaru, jeśli zapisano}

\label{sec:widma-linie}%
Widmo po \SI{100}{\second} czasu życia przedstawiono na \rysref{fig:widmo-4kV}.
Skalę energii pliku programu ProSpect skorygowano liniowo na podstawie
położeń linii Mg~\Ka, Al~\Ka{} i~Si~\Ka{} (\pref{sec:widma-probki}).

\rysunek[0.95\linewidth]{widma-eds/widmo_eds_4kV.pdf}{Widmo EDS zarejestrowane
  przy $\UT=\SI{4}{\kilo\volt}$ i~$\UB=\SI{1}{\kilo\volt}$, czas życia \SI{100}{\second};
  w~powiększeniu zakres pików sumacyjnych \opracowanieWlasne}{fig:widmo-4kV}

W~widmie występują trzy linie charakterystyczne lekkich pierwiastków
(\tabref{tab:widma-linie}). Najsilniejsza jest linia Al~\Ka{} (ok.~$1{,}33\cdot10^{6}$
zliczeń w~kanale \SI{1.49}{\kilo\electronvolt}), a~linie Si~\Ka{} i~Mg~\Ka{} mają
wysokość odpowiednio ok.~\SI{63}{\percent} i~\SI{39}{\percent} piku Al.
Szerokość połówkowa pików wynosi ok.~\SI{87}{\electronvolt}.

\begin{table}[htbp]\centering
  \caption{Linie zidentyfikowane w~widmie z~\rysref{fig:widmo-4kV}}\label{tab:widma-linie}
  \begin{tabular}{@{}S[table-format=1.2]lS[table-format=1.3]l@{}}\toprule
    {$E$ piku [\si{\kilo\electronvolt}]} & Linia & {$E$ tabel. [\si{\kilo\electronvolt}]} & Pochodzenie \\\midrule
    1.26 & Mg~\Ka & 1.254 & domieszka w~próbce Al \\
    1.48 & Al~\Ka & 1.487 & próbka \\
    1.74 & Si~\Ka & 1.740 & struktura krzemowa\,/\,domieszka w~próbce \\
    2.98 & Al~\Ka{} + Al~\Ka & 2.974 & pik sumacyjny \\
    3.22 & Al~\Ka{} + Si~\Ka & 3.227 & pik sumacyjny \\
    3.46 & Si~\Ka{} + Si~\Ka & 3.480 & pik sumacyjny \\\bottomrule
  \end{tabular}
\end{table}

Ciągłe tło promieniowania hamowania zanika przy energii ok.~\SI{4}{\kilo\electronvolt},
co odpowiada granicy Duane'a--Hunta $E_\mathrm{max}=e\UT$ (\pref{sec:rtg}).
Elektrony docierają więc do próbki z~pełną energią ok.~\SI{4}{\kilo\electronvolt},
a~napięcie bramki nie zwiększa energii wiązki. Przy tej energii współczynnik
nadnapięcia wynosi ok.~3,1 dla Mg~K, 2,6 dla Al~K i~2,2 dla Si~K, co zapewnia
skuteczne wzbudzenie tych linii; linii K wapnia i~pierwiastków cięższych
nie można wzbudzić (\pref{subsec:jonizacja}).

\label{sec:widma-stabilnosc}%
W~widmie widoczne są także słabe piki przy \SIlist{2.98;3.22;3.46}{\kilo\electronvolt}, odpowiadające sumom
energii par fotonów Al i~Si zarejestrowanych w~jednym impulsie. Ich obecność
wynika z~bardzo dużej częstości zliczeń, przy której czas martwy przekracza
\SI{50}{\percent}. W~kolejnych pomiarach należy zmniejszyć prąd wiązki lub
zwiększyć odległość detektora od próbki, tak aby czas martwy nie przekraczał
ok.~\SI{30}{\percent}.

Przebieg częstości zliczeń w~czasie pomiaru przedstawiono na \rysref{fig:widma-stabilnosc}.
Częstości wejściowa i~wyjściowa pozostają stałe w~granicach \SI{\pm2}{\percent},
natomiast średnia częstość zliczeń w~kanale Al~\Ka{} maleje w~ciągu \SI{200}{\second}
z~ok.~$15{,}0\cdot10^{3}$ do $13{,}3\cdot10^{3}$~\si{\per\second}, tj.~o~ok.~\SI{11}{\percent}.
Spadek ten może wynikać ze zmian prądu emisji (podrozdz.~\ref{sec:emisja-magnes}), jednak
przy stałej częstości wejściowej bliskiej nasycenia toru akwizycji nie można
tego rozstrzygnąć bez pomiaru $\Ie$.

\rysunek[0.9\linewidth]{widma-eds/stabilnosc_zliczen.pdf}{Częstość zliczeń wejściowa (ICR),
  wyjściowa (OCR) i~średnia częstość w~kanale Al~\Ka{} w~czasie rejestracji widma
  \opracowanieWlasne}{fig:widma-stabilnosc}

Uzyskane widmo potwierdza, że wiązka elektronów z~emitera CNT wzbudza
w~próbce charakterystyczne promieniowanie X, a~linia Al~\Ka{} wskazana
w~celu pracy jest rejestrowana z~dużą liczbą zliczeń już przy
$\UT=\SI{4}{\kilo\volt}$.
.
%  Dane: widma-eds/dane_mca/2026-11-13/ (.mca), skrypt: skrypty_mca/widma_mca.py,
%  położenia pików: piki_13-11.csv, rysunek: widma_materialy.pdf.
% =====================================================================

\subsection{Widma próbek o~różnym składzie}
\label{sec:widma-probki}

Kolejne widma zarejestrowano 13.11.2026 dla czterech próbek: A~(miedź o~nieznanym
składzie i~zanieczyszczonej powierzchni), B~(ołów), C~(złoto, chrom i~miedź) i~D~(aluminium).
Dla próbek A--C wykonano po 6--10 serii o~czasie życia
\SIrange{40}{250}{\second}, które zsumowano. Dla próbki~D wykorzystano serię
o~najmniejszym czasie martwym (\SI{20}{\percent}); pozostałe serie miały czas
martwy do \SI{90}{\percent}, przy którym widmo zniekształcają piki sumacyjne. Napięcie przyspieszające wynosiło
$\UT=\SI{10}{\kilo\volt}$, co potwierdza koniec tła promieniowania hamowania
przy ok.~\SI{10}{\kilo\electronvolt} (granica Duane'a--Hunta), więc
$\Ezero\approx\SI{10}{\kilo\electronvolt}$.

Skala energii zapisana przez program ProSpect (\num{8192} kanały,
\SI{80}{\kilo\electronvolt}) wymagała korekcji: linia Cu~\Ka{} wypadała
w~niej przy \SI{7.89}{\kilo\electronvolt} zamiast \SI{8.04}{\kilo\electronvolt}.
Zależność liniową $E = a\,E_\mathrm{p} + b$ wyznaczono z~położeń pięciu linii
o~jednoznacznej identyfikacji (Al~\Ka, Pb~M$\alpha$, Cr~\Ka, Fe~\Ka, Cu~\Ka).
Otrzymano $a = \num{1.0216}$ i~$b = \SI{-18}{\electronvolt}$, a~odchyłki
punktów kalibracyjnych nie przekraczają \SI{3}{\electronvolt}.
Pozostałe linie po korekcji leżą w~odległości do \SI{12}{\electronvolt}
od wartości tabelarycznych~\cite{goldstein2018sem}.

\rysunek[0.95\linewidth]{widma-eds/widma_materialy.pdf}{Widma EDS próbek A, B i~C
  (pomiary z~13.11.2026, skala logarytmiczna, zliczenia na sekundę czasu życia);
  szare pole -- zakres energii poza czułością toru pomiarowego, linia przerywana --
  położenie linii O~\Ka{} \opracowanieWlasne}{fig:widma-probki}

\begin{table}[htbp]\centering\small
  \caption{Linie zidentyfikowane w~widmach z~13.11.2026}\label{tab:widma-probki}
  \begin{tabular}{@{}lp{0.47\linewidth}p{0.3\linewidth}@{}}\toprule
    Próbka & Linie (energia po korekcji, \si{\kilo\electronvolt}) & Interpretacja \\\midrule
    A & Cu~L$\alpha$ 0,92; Cu~\Ka{} 8,04; Cu~\Kb{} 8,90; Te~L$\alpha$ 3,76,
        L$\beta_1$ 4,04, L$\beta_2$ 4,30, L$\gamma_1$ 4,57; P~\Ka{} 2,01;
        S~\Ka{} 2,30; Cl~\Ka{} 2,62; K~\Ka{} 3,32; Si~\Ka{} 1,73; Al~\Ka{} 1,48
      & miedź; Te i~P (prawdopodobnie dodatki stopowe);
        S, Cl, K, Si, Al z~zanieczyszczeń powierzchni \\
    B & Pb~M$\zeta$ 1,83; Pb~M$\alpha$ 2,35; Pb~M$\gamma$ 2,65 & ołów \\
    C & Au~M$\zeta$ 1,66; Au~M$\alpha\beta$ 2,13; Au~M$\gamma$ 2,41;
        Cr~\Ka{} 5,42; Cr~\Kb{} 5,95; Cu~L$\alpha$ 0,93; Cu~\Ka{} 8,04
      & złoto, chrom, miedź \\
    D & Al~\Ka{} 1,48; Mn~\Ka{} 5,91; Fe~\Ka{} 6,40; pik sumacyjny 2,98
      & stop aluminium z~Mn i~Fe \\\bottomrule
  \end{tabular}
\end{table}

Widma przedstawiono na \rysref{fig:widma-probki}, a~zidentyfikowane linie
zestawiono w~\tabref{tab:widma-probki}. W~próbce~A dominuje linia Cu~L$\alpha$.
Linie Cu~K są słabe, ponieważ przy $\Ezero\approx\SI{10}{\kilo\electronvolt}$
współczynnik nadnapięcia dla krawędzi Cu~K (\SI{8.98}{\kilo\electronvolt})
wynosi tylko ok.~1,1. Grupę pików przy \SIrange{3.76}{4.57}{\kilo\electronvolt}
przypisano serii L telluru: cztery linie zgadzają się z~wartościami tabelarycznymi
w~granicach \SI{9}{\electronvolt}, a~stosunek wysokości L$\beta_1$/L$\alpha$
(ok.~0,4) wyklucza alternatywną interpretację jako Ca~\Ka{} i~Ca~\Kb{}
(\SI{3.69}{} i~\SI{4.01}{\kilo\electronvolt}). Skład próbki nie jest znany. Tellur i~fosfor są typowymi
dodatkami miedzi tellurowej (CuTeP), a~S, Cl, K, Si i~Al pochodzą najpewniej
z~zanieczyszczeń powierzchni.
W~próbce~B widoczne są wyłącznie linie M ołowiu; linie Pb~L nie występują,
ponieważ krawędź Pb~L$_3$ (\SI{13.0}{\kilo\electronvolt}) leży powyżej $\Ezero$.
Widmo próbki~C zawiera linie M złota oraz linie Cr~K i~Cu, co odpowiada
warstwie złota na podłożu miedzianym z~chromem jako warstwą pośrednią.
Częstość zliczeń w~zakresie linii Au~M$\alpha$ w~dziesięciu kolejnych seriach
wynosiła \SIrange{1027}{1096}{\per\second}, tj.~zmieniała się w~granicach
ok.~\SI{\pm3}{\percent}. Widmo próbki~D jest podobne do \rysref{fig:widmo-4kV},
lecz nie zawiera linii Mg i~Si, a~zawiera słabe linie Mn i~Fe, więc jest to
inny stop aluminium.

\label{sec:widma-tlen}%
Osobnego omówienia wymaga zakres niskich energii, w~którym leży linia
tlenu. We wszystkich widmach liczba zliczeń spada poniżej ok.~\SI{0.65}{\kilo\electronvolt}
o~kilka rzędów wielkości, praktycznie do zera
(szare pole na \rysref{fig:widma-probki}). Linia O~\Ka{} (\SI{0.525}{\kilo\electronvolt})
leży poniżej tej granicy, dlatego tlen nie jest rejestrowany, mimo że powierzchnia
każdej z~próbek metalicznych jest pokryta tlenkiem rodzimym (Cu$_2$O/CuO, PbO, Al$_2$O$_3$)
i~zaadsorbowaną wodą. Przy ciśnieniu \SI{1e-5}{\milli\bar} warstwa adsorbatu
odnawia się w~czasie rzędu sekundy, więc pomiar w~komorze nie usuwa tlenu
z~powierzchni. Dolną granicę wyznacza najpewniej absorpcja fotonów w~oknie wejściowym
i~warstwie nieczułej detektora oraz próg dyskryminacji impulsów toru
akwizycji~\cite{goldstein2018sem}; typu okna użytego detektora nie ustalono.

Najsilniejszy pik w~zakresie niskich energii, przy \SI{0.92}{\kilo\electronvolt},
nie pochodzi od tlenu, lecz jest linią Cu~L$\alpha$ (\SI{0.930}{\kilo\electronvolt}).
Potwierdzają to linie Cu~\Ka{} i~Cu~\Kb{} w~tym samym widmie, obecność tego piku
w~próbce~C z~podłożem miedzianym i~jego brak w~próbkach B i~D. Odległość od
linii O~\Ka{} wynosi ok.~\SI{400}{\electronvolt}, tj.~pięciokrotność szerokości
połówkowej piku (ok.~\SI{80}{\electronvolt}), więc nakładanie obu linii jest
wykluczone. Pik leży tuż nad granicą czułości, przez co jego wysokość względem
linii o~większych energiach jest zaniżona.

Do oznaczania tlenu, najczęstszego pierwiastka skał i~regolitu, analizator
docelowy wymaga detektora z~cienkim oknem polimerowym lub bez okna.
